\subsection{Semantic Subset Analysis}
\label{sec:semantic-subset}

The cross-dataset entropy validation (Section~\ref{sec:cross-dataset-entropy}) reveals a scale-dependent trade-off: entropy-only guidance excels on 10-class datasets, while the multi-factor \cags{} metric wins on the 100-class ImageNet-100. A natural question follows: is this trade-off driven by the \emph{number} of classes or by the \emph{semantic composition} of the dataset? To disentangle these factors, we partition the 100 ImageNet-100 classes into two semantically coherent subsets---44 animal classes (``animals'') and 56 non-animal classes (``objects'')---and evaluate all three guidance configurations on each subset independently. Both subsets share the same generation pipeline, guidance range $[0.0, 0.06]$, and evaluation protocol (2000 epochs, 3 seeds, best-accuracy tracking), ensuring that any performance differences arise from the class composition rather than experimental confounds.

\Cref{tab:semantic-subset} presents the results. The findings confirm that the multi-factor \cags{} advantage is driven by class scale rather than semantic homogeneity: \cags{} outperforms entropy-only on both the 44-class animals subset and the 56-class objects subset, consistent with its superiority on the full 100-class dataset.

\begin{table}[t]
\centering
\caption{Semantic subset analysis on ImageNet-100 partitions. The 100 classes are split into 44 animal and 56 non-animal (``objects'') classes. \cags{} outperforms entropy-only on both subsets, confirming that the multi-factor advantage scales with class count rather than being tied to semantic diversity. All results use 3 seeds with best-accuracy tracking.}
\label{tab:semantic-subset}
\begin{tabular}{llccc}
\toprule
\textbf{Subset} & \textbf{Method} & \textbf{Top-1 (\%)} & \textbf{Top-5 (\%)} & \textbf{$\Delta$ vs.\ unguided} \\
\midrule
\multirow{3}{*}{Animals (44-cls)} & Unguided & $23.64 \pm 0.37$ & $49.02 \pm 0.85$ & --- \\
 & \cags{} (multi-factor) & $\mathbf{25.94 \pm 0.46}$ & $\mathbf{51.24 \pm 0.49}$ & $\mathbf{+2.30}$ \\
 & Entropy-only & $25.29 \pm 0.39$ & $51.05 \pm 0.10$ & $+1.65$ \\
\midrule
\multirow{3}{*}{Objects (56-cls)} & Unguided & $25.08 \pm 0.29$ & $46.45 \pm 0.50$ & --- \\
 & \cags{} (multi-factor) & $\mathbf{26.85 \pm 0.25}$ & $\mathbf{50.56 \pm 0.50}$ & $\mathbf{+1.77}$ \\
 & Entropy-only & $26.42 \pm 0.12$ & $49.60 \pm 0.40$ & $+1.34$ \\
\bottomrule
\end{tabular}
\end{table}

\paragraph{\cags{} wins on both semantic subsets.}
On the 44-class animals subset, \cags{} achieves $25.94 \pm 0.46\%$, outperforming entropy-only ($25.29 \pm 0.39\%$) by $0.65\%$ absolute ($2.6\%$ relative) and unguided ($23.64 \pm 0.37\%$) by $2.30\%$ absolute ($9.7\%$ relative). On the 56-class objects subset, \cags{} achieves $26.85 \pm 0.25\%$, outperforming entropy-only ($26.42 \pm 0.12\%$) by $0.43\%$ absolute ($1.6\%$ relative) and unguided ($25.08 \pm 0.29\%$) by $1.77\%$ absolute ($7.1\%$ relative). These results demonstrate that the multi-factor advantage of \cags{} is not an artifact of the specific 100-class composition but rather a function of the number of classes: at 44 and 56 classes, the additional complexity factors---mode count, intra-class variance, and inter-class separability---already provide complementary signals that entropy alone cannot capture.

\paragraph{Contrast with 10-class datasets.}
On the 10-class ImageNette and ImageWoof datasets (Section~\ref{sec:cross-dataset-entropy}), entropy-only guidance outperforms \cags{} by $1.4$--$1.6\%$. The semantic subset results bridge this gap: at 44 classes, \cags{} surpasses entropy-only by $0.65\%$, and at 56 classes by $0.43\%$, suggesting that the transition from entropy-only to multi-factor superiority occurs somewhere between 10 and 44 classes. This threshold is consistent with the theoretical motivation (Proposition~\ref{prop:entropy}): with fewer classes, the complexity factors exhibit limited variation and add noise, but as the class count grows, the factors capture genuinely different aspects of class complexity that entropy alone cannot represent.

\paragraph{Implications for practical deployment.}
The combined findings from the cross-dataset validation and semantic subset analysis yield a clear practical recommendation. For small-scale datasets ($\leq$10 classes), entropy-only guidance is both simpler and more effective. For medium-to-large-scale datasets ($\geq$44 classes), the full multi-factor \cags{} metric provides consistent advantages across different semantic compositions. Both approaches remain training-free and require only a one-time complexity analysis, preserving the practical appeal of the \ags{} framework.
